% p343_1; source PDF page 343; manually verified editable table.
\begin{tikzpicture}[x=1cm,y=-1cm,inner sep=0pt,outer sep=0pt]
\path[use as bounding box] (0,0) rectangle (16.000,8.075);
\draw[line width=.4pt] (0.000,0.000) rectangle (2.000,0.815);
\node[align=center,text width=1.740cm,font={\linespread{1}\fontsize{10}{13.50}\selectfont\bfseries}] at (1.000,0.407) {构成要件};
\draw[line width=.4pt] (2.000,0.000) rectangle (6.400,0.815);
\node[align=center,text width=4.140cm,font={\linespread{1}\fontsize{10}{13.50}\selectfont\bfseries}] at (4.200,0.407) {具体内容};
\draw[line width=.4pt] (6.400,0.000) rectangle (16.000,0.815);
\node[align=center,text width=9.340cm,font={\linespread{1}\fontsize{10}{13.50}\selectfont\bfseries}] at (11.200,0.407) {要点};
\draw[line width=.4pt] (0.000,0.815) rectangle (2.000,3.330);
\node[align=center,text width=1.740cm,font={\linespread{1}\fontsize{10}{13.50}\selectfont}] at (1.000,2.072) {行为主体};
\draw[line width=.4pt] (2.000,0.815) rectangle (6.400,3.330);
\node[align=left,text width=4.140cm,font={\linespread{1}\fontsize{10}{13.50}\selectfont}] at (4.200,2.072) {不限于贷款人，包括担保人};
\draw[line width=.4pt] (6.400,0.815) rectangle (16.000,3.330);
\node[align=left,text width=9.340cm,font={\linespread{1}\fontsize{10}{13.50}\selectfont}] at (11.200,2.072) {甲没有非法占有目的，为了取得银行贷款，请乙提供担保。乙提供了虚假的担保材料，甲不知情，银行因为受骗而发放贷款。甲因客观原因无法还款，银行遭受重大损失。乙独立构成本罪};
\draw[line width=.4pt] (0.000,3.330) rectangle (2.000,4.745);
\node[align=center,text width=1.740cm,font={\linespread{1}\fontsize{10}{13.50}\selectfont}] at (1.000,4.037) {行为方式};
\draw[line width=.4pt] (2.000,3.330) rectangle (6.400,4.745);
\node[align=left,text width=4.140cm,font={\linespread{1}\fontsize{10}{13.50}\selectfont}] at (4.200,4.037) {以欺骗手段取得银行贷款};
\draw[line width=.4pt] (6.400,3.330) rectangle (16.000,4.745);
\node[align=left,text width=9.340cm,font={\linespread{1}\fontsize{10}{13.50}\selectfont}] at (11.200,4.037) {如果金融机构工作人员知道真相仍然放贷，则行为人不构成本罪，相反，金融机构工作人员应构成违法发放贷款罪};
\draw[line width=.4pt] (0.000,4.745) rectangle (2.000,6.660);
\node[align=center,text width=1.740cm,font={\linespread{1}\fontsize{10}{13.50}\selectfont}] at (1.000,5.703) {危害结果};
\draw[line width=.4pt] (2.000,4.745) rectangle (6.400,6.660);
\node[align=left,text width=4.140cm,font={\linespread{1}\fontsize{10}{13.50}\selectfont}] at (4.200,5.703) {要求给银行造成重大损失};
\draw[line width=.4pt] (6.400,4.745) rectangle (16.000,6.660);
\node[align=left,text width=9.340cm,font={\linespread{1}\fontsize{10}{13.50}\selectfont}] at (11.200,5.703) {该罪的成立条件原来包括两种情形，或者给银行造成重大损失，或者有严重情节。《刑法修正案（十一）》删除“有严重情节”。本罪的成立条件只剩下“给银行造成重大损失”};
\draw[line width=.4pt] (0.000,6.660) rectangle (2.000,8.075);
\node[align=center,text width=1.740cm,font={\linespread{1}\fontsize{10}{13.50}\selectfont}] at (1.000,7.367) {主观要件};
\draw[line width=.4pt] (2.000,6.660) rectangle (6.400,8.075);
\node[align=left,text width=4.140cm,font={\linespread{1}\fontsize{10}{13.50}\selectfont}] at (4.200,7.367) {对贷款没有非法占有目的，也即具有归还意思};
\draw[line width=.4pt] (6.400,6.660) rectangle (16.000,8.075);
\node[align=left,text width=9.340cm,font={\linespread{1}\fontsize{10}{13.50}\selectfont}] at (11.200,7.367) {本罪与贷款诈骗罪的区别在于：贷款诈骗罪具有非法占有目的};
\end{tikzpicture}
